\documentclass[11pt,a4paper]{article}
\usepackage[margin=17mm,top=16mm,bottom=18mm]{geometry}
\usepackage{fontspec}
\setmainfont{texgyrepagella}[Extension=.otf,UprightFont=*-regular,BoldFont=*-bold,ItalicFont=*-italic,BoldItalicFont=*-bolditalic]
\usepackage{enumitem}
\usepackage{xcolor}
\definecolor{accent}{HTML}{235A73}
\definecolor{muted}{HTML}{555555}
\usepackage{titlesec}
\usepackage{fancyhdr}
\usepackage{hyperref}
\hypersetup{colorlinks=true,urlcolor=accent,pdftitle={Jeffrey Hu — Curriculum Vitae},pdfauthor={Jeffrey Hu}}
\setlength{\parindent}{0pt}
\setlength{\emergencystretch}{1em}
\hyphenpenalty=2000
\setlength{\parskip}{3pt}
\setlist[itemize]{leftmargin=13pt,itemsep=2pt,topsep=3pt,parsep=0pt}
\setlist[itemize,2]{label=\textendash,leftmargin=12pt,itemsep=1pt,topsep=1pt}
\titleformat{\section}{\large\scshape\color{accent}}{}{0pt}{}[\vspace{-4pt}\color{accent}\titlerule]
\titlespacing*{\section}{0pt}{9pt}{5pt}
\pagestyle{fancy}
\fancyhf{}
\fancyfoot[L]{\footnotesize\color{muted}Jeffrey Hu}
\fancyfoot[R]{\footnotesize\color{muted}\thepage}
\renewcommand{\headrulewidth}{0pt}
\renewcommand{\footrulewidth}{0pt}
\setlength{\footskip}{22pt}
\newcommand{\entry}[4]{%
  \textbf{#1}\hfill #2\\
  \textit{#3}\hfill\textcolor{muted}{#4}\par
}
\newcommand{\publication}[4]{%
  {\raggedright\hyphenpenalty=10000\href{#1}{\textbf{#2}}\par}
  #3\\
  \textit{#4}\par\vspace{4pt}
}
\begin{document}

{\fontsize{27}{30}\selectfont\textbf{Jeffrey Hu}}\hfill
\href{https://jefequien.github.io/}{jefequien.github.io}\\[3pt]
\href{mailto:hujh14@gmail.com}{hujh14@gmail.com}\enspace\textcolor{muted}{/}\enspace
\href{https://github.com/jefequien}{github.com/jefequien}\enspace\textcolor{muted}{/}\enspace
\href{https://www.linkedin.com/in/jeffrey-h-hu/}{linkedin.com/in/jeffrey-h-hu}

\section{Education}
\entry{University of Cambridge}{Current}{PhD student}{Cambridge, UK}
Supervisor: Professor Ayush Tewari. Research in 3D computer vision, computer graphics,
differentiable rendering, and generative models of the real world.

\vspace{3pt}
\entry{Massachusetts Institute of Technology}{2018, 2019}
{Computer Science and Engineering}{Cambridge, MA}
Master of Science, 2019; Bachelor of Science, 2018. GPA: 4.7/5.0.\\
Research advisor: Professor Antonio Torralba.

\section{Experience}
\entry{Inria — GraphDeco}{Previously}{Research Engineer}{}
\begin{itemize}
  \item Worked with Professor George Drettakis on 3D Gaussian splatting and differentiable
  rendering, including intrinsic decomposition, editable reflections, and reconstruction
  of scenes with oscillatory motion.
\end{itemize}

\entry{gsplat}{From Apr 2024}{Open-source Contributor}{}
\begin{itemize}
  \item Contributed to an open-source Gaussian splatting library for 3D reconstruction and
  novel-view synthesis, including:
  \begin{itemize}
    \item \href{https://github.com/nerfstudio-project/gsplat/pull/238}{3DGS-MCMC}
    for improved densification and
    \href{https://github.com/nerfstudio-project/gsplat/pull/291}{bilateral grids}
    for handling exposure changes.
    \item \href{https://github.com/nerfstudio-project/gsplat/pull/398}{Fisheye camera support}
    and \href{https://github.com/nerfstudio-project/gsplat/pull/273}{2.5D Gaussian splatting}
    for surface reconstruction.
    \item Combined MCMC, Self-Organizing Gaussian Splats, and spherical-harmonics
    k-means clustering for
    \href{https://github.com/nerfstudio-project/gsplat/pull/309}{Gaussian splat compression}.
    \item \href{https://github.com/nerfstudio-project/gsplat/pull/469}{Camera defocus blur modeling}.
  \end{itemize}
\end{itemize}

\entry{Parallel Systems}{Mar 2022--Apr 2024}
{Senior Software Engineer, Perception}{Los Angeles, CA}
\begin{itemize}
  \item Led development of a real-time multitask computer vision network and data engine
  for an autonomous train perception system.
  \item Iteratively built a multitask network for 3D object detection, 3D track prediction,
  depth prediction, and bird's-eye-view segmentation.
  \begin{itemize}
    \item Wrote a PyTorch library for training models with various backbones,
    task-specific heads, and loss functions.
    \item Developed a representation for railroad track geometry by regressing 3D splines
    and classifying switch states.
    \item Shipped multiple model iterations onto a Jetson using Rust, TensorRT, and DeepStream.
  \end{itemize}
  \item Developed a 3D reconstruction and autolabeling pipeline for scalable ground-truth
  generation using Dagster and Kubernetes.
  \begin{itemize}
    \item Used zero-shot, open-vocabulary foundation models for object detection,
    segmentation, tracking, and monocular depth estimation.
    \item Aggregated predictions across time to reconstruct 3D scene representations
    as high-quality labels for training and evaluating the real-time network.
    \item Built infrastructure to balance datasets, mine scenarios, and automatically
    generate test reports to support the regulatory safety case over time.
  \end{itemize}
  \item Contributed to sensor calibration, onboard data collection, triggering,
  and upload infrastructure.
\end{itemize}

\newpage
{\large\textbf{Jeffrey Hu}}\hfill\textcolor{muted}{Curriculum vitae}

\section{Experience, continued}
\entry{TuSimple}{Jan 2020--Mar 2022}{Research Engineer II}{San Diego, CA}
\begin{itemize}
  \item Developed a method for correcting camera drift in visual odometry by training
  a network to regress the difference between the current image and a rendered synthetic map view.
  \item Wrote an onboard C++ library to efficiently manage 6DoF transformation matrices
  between coordinate frames of a scene graph in an asynchronous system.
\end{itemize}

\entry{MIT Computer Vision Group}{Sep 2017--Sep 2019}{Research Assistant}{Cambridge, MA}
\begin{itemize}
  \item Investigated scalable semi-automatic methods for labeling large segmentation datasets.
  \item Built a browser-based annotation tool for rapid segmentation editing.
\end{itemize}

\entry{Amazon Prime Air}{Jun 2016--Aug 2016}{Software Engineer Intern}{Seattle, WA}

\section{Publications}
\publication{https://arxiv.org/abs/2606.31637}
{Intrinsic Decomposition and Editing of 3D Gaussian Splats}
{A. Lanvin, \textbf{J. Hu}, S. Lucas, A. Bousseau, and G. Drettakis}
{I3D / Proceedings of the ACM on Computer Graphics and Interactive Techniques, 2026.}

\publication{https://repo-sam.inria.fr/nerphys/adaptive-spatio-temporal/AdaptiveSpatioTemporal_authors.pdf}
{Adaptive Spatio-Temporal 3D Gaussian Splatting for Scenes with Oscillatory Motion}
{P. Tzathas, \textbf{J. Hu}, A. Meuleman, G. Cordonnier, and G. Drettakis}
{Eurographics / Computer Graphics Forum, 2026.}

\publication{https://arxiv.org/abs/2601.09697}
{Efficient Camera-Controlled Video Generation of Static Scenes via Sparse Diffusion and 3D Rendering}
{J. Chen, \textbf{J. Hu}, J. Lasenby, and A. Tewari}
{arXiv preprint, 2026.}

\publication{https://arxiv.org/abs/2606.30861}
{Editable Physically-based Reflections in Raytraced Gaussian Radiance Fields}
{Y. Poirier-Ginter, \textbf{J. Hu}, J.-F. Lalonde, and G. Drettakis}
{SIGGRAPH Asia, 2025.}

\publication{https://arxiv.org/abs/2409.06765}
{gsplat: An Open-Source Library for Gaussian Splatting}
{V. Ye, R. Li, J. Kerr, M. Turkulainen, B. Yi, Z. Pan, O. Seiskari, J. Ye,
\textbf{J. Hu}, M. Tancik, and A. Kanazawa}
{Journal of Machine Learning Research, 2025.}

\section{Leadership \& interests}
\entry{MIT Varsity Squash}{Sep 2014--Jun 2018}{Captain}{Cambridge, MA}
\textbf{Programming:} Python, PyTorch, C++, Rust, CUDA, Git.\\
\textbf{Interests:} Squash, cycling, skiing, hiking.

\end{document}
